\documentclass[10pt, xcolor={dvipsnames}]{beamer}

% --- Cấu hình Tiếng Việt và Giao diện ---
\usepackage[utf8]{vietnam}
\usepackage{amsmath, amssymb}
\usepackage{listings}
\usepackage{booktabs}

\usetheme{Madrid}
\usecolortheme{default} % Sử dụng tông trắng chủ đạo
\setbeamertemplate{items}[square]
\setbeamertemplate{navigation symbols}{} 

% Tùy chỉnh màu sắc tiêu đề nhẹ nhàng (Xanh Navy)
\setbeamercolor{structure}{fg=MidnightBlue}
\setbeamercolor{title}{fg=white, bg=MidnightBlue}

% --- Cấu hình hiển thị Code (Nền trắng, sạch sẽ) ---
\lstset{
    language=C++,
    backgroundcolor=\color{white},   
    commentstyle=\color{gray},
    keywordstyle=\color{blue}\bfseries,
    numberstyle=\tiny\color{gray},
    stringstyle=\color{OliveGreen},
    basicstyle=\ttfamily\tiny,
    breaklines=true,                 
    captionpos=b,                    
    frame=lines, % Chỉ kẻ dòng trên và dưới cho thanh lịch
    numbers=left,
    tabsize=2
}

% --- Thông tin bài thuyết trình ---
\title[Shortest Path]{THUẬT TOÁN ĐƯỜNG ĐI NGẮN NHẤT}
\subtitle{Phân tích giải thuật Dijkstra \& Floyd-Warshall}
\author[Phạm Gia Kiệt]{Thuyết trình bởi: \textbf{Phạm Gia Kiệt}}
\institute[Lớp 11 Tin]{Khối 11 - Chuyên Tin}
\date{\today}

\begin{document}

% Slide Tiêu đề
\begin{frame}
    \titlepage
\end{frame}

% Slide 1: Đặt vấn đề
\begin{frame}{1. Bài toán tìm đường đi ngắn nhất}
    \begin{itemize}
        \item \textbf{Khái niệm:} Tìm lộ trình giữa các cặp đỉnh sao cho tổng trọng số (chi phí, khoảng cách) là tối ưu nhất.
        \item \textbf{Phân loại theo mục tiêu:}
        \begin{enumerate}
            \item \textbf{Single-Source:} Dijkstra, Bellman-Ford.
            \item \textbf{All-Pairs:} Floyd-Warshall.
        \end{enumerate}
    \end{itemize}
    \vspace{0.3cm}
    \begin{block}{Ứng dụng thực tế}
        Tìm đường trên bản đồ số, tối ưu hóa băng thông mạng thông tin, bài toán vận chuyển hàng hóa.
    \end{block}
\end{frame}

% Slide 2: Dijkstra
\begin{frame}{2. Thuật toán Dijkstra}
    \textbf{Đặc điểm:} Thuật toán tham lam (Greedy). Luôn chọn đỉnh có khoảng cách tạm thời nhỏ nhất để tối ưu.
    
    \begin{itemize}
        \item \textbf{Yêu cầu:} Trọng số trên các cạnh \textbf{không được âm}.
        \item \textbf{Độ phức tạp:} $O(E \log V)$ khi dùng \texttt{priority\_queue}.
    \end{itemize}

    \begin{exampleblock}{Công thức thư giãn (Relaxation)}
        Dựa trên bất đẳng thức tam giác:
        $$d[v] = \min(d[v], d[u] + w(u, v))$$
    \end{exampleblock}
\end{frame}

% Slide 3: Code Dijkstra
\begin{frame}[fragile]{Dijkstra - C++ Implementation}
\begin{lstlisting}
typedef pair<int, int> pii;
const int INF = 1e9 + 7;
vector<pii> adj[100005];
int d[100005];

void dijkstra(int s) {
    fill(d + 1, d + n + 1, INF);
    priority_queue<pii, vector<pii>, greater<pii>> pq;
    d[s] = 0; pq.push({0, s});

    while (!pq.empty()) {
        int u = pq.top().second, du = pq.top().first;
        pq.pop();
        if (du > d[u]) continue;
        for (auto edge : adj[u]) {
            int v = edge.first, w = edge.second;
            if (d[v] > d[u] + w) {
                d[v] = d[u] + w;
                pq.push({d[v], v});
            }
        }
    }
}
\end{lstlisting}
\end{frame}

% Slide 4: Floyd-Warshall
\begin{frame}{3. Thuật toán Floyd-Warshall}
    \textbf{Nguyên lý:} Quy hoạch động (Dynamic Programming). Xét lần lượt từng đỉnh $k \in \{1 \dots N\}$ làm điểm trung gian giữa mọi cặp $(i, j)$.
    
    \begin{itemize}
        \item \textbf{Đặc điểm:} Tìm đường đi ngắn nhất giữa \textbf{mọi cặp đỉnh}.
        \item \textbf{Độ phức tạp:} $O(V^3)$.
        \item \textbf{Ưu điểm:} Xử lý được cạnh âm, cài đặt cực kỳ đơn giản.
    \end{itemize}

    \begin{alertblock}{Công thức chuyển trạng thái}
        $$dist[i][j] = \min(dist[i][j], dist[i][k] + dist[k][j])$$
    \end{alertblock}
\end{frame}

% Slide 5: Code Floyd
\begin{frame}[fragile]{Floyd-Warshall - Implementation}
\begin{lstlisting}
// dist[i][j] khoi tao bang trong so canh
// dist[i][i] = 0, cac cap ko co canh = INF
void floyd() {
    for (int k = 1; k <= n; k++) {
        for (int i = 1; i <= n; i++) {
            for (int j = 1; j <= n; j++) {
                if (dist[i][k] < INF && dist[k][j] < INF)
                    dist[i][j] = min(dist[i][j], dist[i][k] + dist[k][j]);
            }
        }
    }
}
\end{lstlisting}
\end{frame}

% Slide 6: So sánh
\begin{frame}{4. So sánh tổng quát}
    \centering
    \begin{table}
        \begin{tabular}{lcc}
            \toprule
            \textbf{Đặc điểm} & \textbf{Dijkstra} & \textbf{Floyd-Warshall} \\
            \midrule
            Mô hình & Nguồn đơn (SSSP) & Mọi cặp đỉnh (APSP) \\
            Trọng số âm & Không hỗ trợ & Có hỗ trợ \\
            Độ phức tạp & $O(E \log V)$ & $O(V^3)$ \\
            Cài đặt & Trung bình & Rất dễ \\
            \bottomrule
        \end{tabular}
    \end{table}
\end{frame}

% Slide Kết thúc
\begin{frame}
    \centering
    \Huge \textcolor{MidnightBlue}{\textbf{CẢM ƠN CÁC BẠN!}} \\
    \vspace{0.5cm}
    \large \textbf{Q\&A}
\end{frame}

\end{document}
